\chapter{Example Two-Level Atom System} % Main chapter title

\label{chapter_rabi_oscillations} % For referencing this chapter elsewhere, use \autoref{chapter_rabi_oscillations}

%----------------------------------------------------------------------------------------
%	SECTION: Two-Level Atom Master Equation in the Rotating Frame with RWA
%----------------------------------------------------------------------------------------

\section{Master Equation in the Rotating Frame with RWA}

Rabi oscillations describe the coherent oscillatory behavior of a two-level quantum system interacting with a resonant electromagnetic field.
This phenomenon is fundamental in quantum mechanics and quantum optics.


\subsection{Hamiltonian in the Laboratory Frame}
Consider a two-level system with states \(\ket{g}\) (ground state) and \(\ket{e}\) (excited state).
The energy separation between the two states is given by:
\begin{equation}
	\label{eq:EnergySeparation}
\end{equation}
where \(E_e\) and \(E_g\) are the energies of the excited and ground states, respectively.

\subsubsection{Atomic Hamiltonian}
In the convention used in Paper~I and in the implementation, the ground-state energy is set to zero and the free Hamiltonian is written as
\begin{equation}
	H_0 = \omega_{eg}\hat{N}, \qquad \hat{N} = \ket{e}\bra{e} = \sigma_+\sigma_-.
	\label{eq:AtomicHamiltonian}
\end{equation}
This differs from the centered form \(H_0 = \omega_{eg}\sigma_z/2\) only by an additive multiple of the identity and therefore leads to the same equations of motion.


\subsubsection{Electric Field}
The code distinguishes between the laboratory-frame positive-frequency field \(\varepsilon^{(+)}(t)\) and its slowly varying envelope \(E(t)\):
\begin{equation}
	\varepsilon^{(+)}(t) = e^{-i\omega_{L} t} E(t), \qquad
	E(t) = E_0 e^{-i\phi}, \qquad
	E_{\mathrm{phys}}(t) = \varepsilon^{(+)}(t) + \varepsilon^{(-)}(t).
	\label{eq:ElectricFieldComplex}
\end{equation}
Here \(E_0\) is the field amplitude, \(\phi\) is the pulse phase, and \(\omega_{L}\) is the laser carrier frequency.
This sign convention matches Paper~I and the codebase.

\subsubsection{Interaction Hamiltonian}
Using the electric dipole interaction with raising and lowering operators \(\sigma_\pm = (\sigma_x \pm i\sigma_y)/2\), the dipole operator can be written as
\begin{equation}
	\hat{\mu} = \mu_{eg}\sigma_+ + \mu_{eg}\sigma_-,
	\label{eq:InteractionHamiltonianDipole}
\end{equation}
where, for a two-level system, the phase of the basis states can be chosen such that \(\mu_{eg}\) is real.

The full laboratory-frame interaction Hamiltonian is then
\begin{equation}
	H_{\mathrm{int}}(t) = -\hat{\mu}\left(\varepsilon^{(+)}(t) + \varepsilon^{(-)}(t)\right).
	\label{eq:asdasdasd}
\end{equation}
Applying the rotating-wave approximation in the laboratory frame keeps only the resonant terms,
\begin{equation}
	H_{\mathrm{int}}^{\mathrm{RWA}}(t) = -\mu_{eg}\left(\varepsilon^{(+)}(t)\sigma_+ + \varepsilon^{(-)}(t)\sigma_-\right).
\end{equation}

\subsubsection{Total Hamiltonian}
The total Hamiltonian of the system is:
\begin{equation}
	H(t) = H_0 + H_{\mathrm{int}}(t).
	\label{eq:TotalHamiltonian}
\end{equation}
The time evolution of the system is governed by the time-dependent Schr\"odinger equation \autoref{eq:oqs_schroedinger_equation}.

\subsection{Transform to the Rotating Frame}

Define the rotating-frame unitary using the excitation-number operator:
\begin{equation}
	U(t) = \exp\left(i\omega_{L} t \hat{N}\right).
	\label{eq:UnitaryRotatingFrame}
\end{equation}
Transform operators as \(\tilde{X} = UXU^\dagger\), and the density matrix as \(\tilde{\rho} = U\rho U^\dagger\).

Useful identities for the transformation are
\begin{equation}
	U\sigma_\pm U^\dagger = e^{\pm i\omega_{L} t}\sigma_\pm, \qquad
	U\hat{N}U^\dagger = \hat{N}.
	\label{eq:TransformationIdentities}
\end{equation}

\subsubsection{Transformed Hamiltonian Terms}

\begin{itemize}
	\item \textbf{Free Hamiltonian:}
	      \begin{equation}
		      \tilde{H}_0 = (\omega_{eg} - \omega_{L})\hat{N} = \Delta_L \hat{N},
		      \label{eq:TransformedFreeHamiltonian}
	      \end{equation}
	      where \(\Delta_L = \omega_{eg} - \omega_{L}\) is the laser detuning.

	\item \textbf{Transformed interaction Hamiltonian:}
	      \begin{equation}
		      \tilde{H}_{\mathrm{int}} = -\mu_{eg}\left(E(t)\sigma_+ + E^*(t)\sigma_-\right).
		      \label{eq:TransformedInteractionHamiltonian}
	      \end{equation}

	\item Defining the complex drive amplitude
	      \begin{equation}
		      \Omega_c(t) = \mu_{eg}E(t),
		      \label{eq:TimeDepRabiFreq}
	      \end{equation}
	      the total rotating-frame Hamiltonian becomes
	      \begin{equation}
		      \boxed{
			      \tilde{H} = \Delta_L \hat{N} - \Omega_c(t)\sigma_+ - \Omega_c^*(t)\sigma_-
		      }
		      \label{eq:RotatingFrameHamiltonian}
	      \end{equation}
\end{itemize}

\subsection{Master Equation in the Rotating Frame (with RWA)}

From here on, \(\rho\) denotes the rotating-frame density matrix.
The master equation is
\begin{equation}
	\frac{\mathrm{d}}{\mathrm{d}t}\rho = -i[\tilde{H}, \rho] + \mathcal{D}(\rho),
	\label{eq:MasterEquationRotatingFrame}
\end{equation}
with dissipator
\begin{equation}
	\mathcal{D}(\rho) =
	\gamma_{\downarrow}\left(\sigma_-\rho\sigma_+ - \frac{1}{2}\{\sigma_+\sigma_-, \rho\}\right)
	+ \gamma_{\uparrow}\left(\sigma_+\rho\sigma_- - \frac{1}{2}\{\sigma_-\sigma_+, \rho\}\right)
	+ \frac{\gamma_\varphi}{2}\left(\sigma_z \rho \sigma_z - \rho\right).
	\label{eq:DissipatorOperator}
\end{equation}
Define the total decoherence rate:
\begin{equation}
	\Gamma = \frac{\gamma_{\downarrow} + \gamma_{\uparrow}}{2} + \gamma_\varphi.
	\label{eq:TotalDecoherenceRate}
\end{equation}
For the Paper-I benchmark one sets \(\gamma_{\uparrow} = 0\).

\subsection{Equations of Motion}

Write the density matrix as
\begin{equation}
	\rho =
	\begin{pmatrix}
		\rho_{gg} & \rho_{ge} \\
		\rho_{eg} & \rho_{ee}
	\end{pmatrix},
	\qquad
	\rho_{ge} = \rho_{eg}^*.
	\label{eq:oqs_DensityMatrixElements}
\end{equation}
Following Paper~I, it is convenient to write the laboratory-frame optical coherences as
\begin{equation}
	\rho_{eg}^{\mathrm{lab}}(t) = \sigma_{eg}(t)e^{-i\omega_L t}, \qquad
	\rho_{ge}^{\mathrm{lab}}(t) = \sigma_{ge}(t)e^{+i\omega_L t}.
\end{equation}
In the rotating frame used above, \(\sigma_{eg} = \rho_{eg}\) and \(\sigma_{ge} = \rho_{ge}\).
The resulting equations of motion are
\begin{equation}
	\boxed{
		\begin{aligned}
			\frac{\mathrm{d}}{\mathrm{d}t}\sigma_{eg} & = -(\Gamma + i\Delta_L)\sigma_{eg}
			+ i\Omega_c(t)(\rho_{gg} - \rho_{ee}) \\[6pt]
			\frac{\mathrm{d}}{\mathrm{d}t}\sigma_{ge} & = -(\Gamma - i\Delta_L)\sigma_{ge}
			- i\Omega_c^*(t)(\rho_{gg} - \rho_{ee}) \\[6pt]
			\frac{\mathrm{d}}{\mathrm{d}t}\rho_{ee} & = -\gamma_{\downarrow}\rho_{ee}
			+ \gamma_{\uparrow}\rho_{gg}
			+ i\Omega_c(t)\sigma_{ge}
			- i\Omega_c^*(t)\sigma_{eg} \\[6pt]
			\frac{\mathrm{d}}{\mathrm{d}t}\rho_{gg} & = -\gamma_{\uparrow}\rho_{gg}
			+ \gamma_{\downarrow}\rho_{ee}
			- i\Omega_c(t)\sigma_{ge}
			+ i\Omega_c^*(t)\sigma_{eg}
		\end{aligned}
	}
	\label{eq:BlochEquations}
\end{equation}
For \(\gamma_{\uparrow} = 0\), these reduce to the Paper-I equations, with the detuning term added explicitly.

\subsection{Liouvillian Matrix Form}

The custom implementation uses QuTiP's column-stacking convention.
In the basis \(\{\ket{g}, \ket{e}\}\), this means
\begin{equation}
	\boldsymbol{\rho} =
	\begin{pmatrix}
		\rho_{gg} \\
		\sigma_{eg} \\
		\sigma_{ge} \\
		\rho_{ee}
	\end{pmatrix}.
	\label{eq:oqs_DensityMatrixVector}
\end{equation}
Then the master equation becomes
\begin{equation}
	\dot{\boldsymbol{\rho}} = \left(L_0 + E(t)L_+ + E^*(t)L_-\right)\boldsymbol{\rho},
	\label{eq:LiouvillianMatrixForm}
\end{equation}
with
\begin{equation}
	L_0 =
	\begin{pmatrix}
		-\gamma_{\uparrow} & 0 & 0 & \gamma_{\downarrow} \\
		0 & -(\Gamma + i\Delta_L) & 0 & 0 \\
		0 & 0 & -(\Gamma - i\Delta_L) & 0 \\
		\gamma_{\uparrow} & 0 & 0 & -\gamma_{\downarrow}
	\end{pmatrix},
	\qquad
	L_+ =
	\begin{pmatrix}
		0 & 0 & -i\mu_{eg} & 0 \\
		i\mu_{eg} & 0 & 0 & -i\mu_{eg} \\
		0 & 0 & 0 & 0 \\
		0 & 0 & i\mu_{eg} & 0
	\end{pmatrix},
\end{equation}
\begin{equation}
	L_- =
	\begin{pmatrix}
		0 & i\mu_{eg} & 0 & 0 \\
		0 & 0 & 0 & 0 \\
		-i\mu_{eg} & 0 & 0 & i\mu_{eg} \\
		0 & -i\mu_{eg} & 0 & 0
	\end{pmatrix}.
\end{equation}
This is exactly the \(L_0 + E(t)L_+ + E^*(t)L_-\) structure used in the \texttt{paper\_eqs} implementation.


\subsection{Constant-Drive Limit}
For a constant envelope \(E(t) = E_0 e^{-i\phi}\), the complex drive amplitude \(\Omega_c = \mu_{eg}E_0 e^{-i\phi}\) is time independent.
In the basis \(\{\ket{g}, \ket{e}\}\), the rotating-frame Hamiltonian is then
\begin{equation}
	\tilde{H}_{\mathrm{RWA}} =
	\begin{pmatrix}
		0 & -\Omega_c^* \\
		-\Omega_c & \Delta_L
	\end{pmatrix}.
	\label{eq:RWAHamiltonianMatrix}
\end{equation}
Subtracting \((\Delta_L/2)\mathbb{1}\) gives the equivalent centered form
\begin{equation}
	\tilde{H}_{\mathrm{RWA}} - \frac{\Delta_L}{2}\mathbb{1}
	=
	\frac{1}{2}
	\begin{pmatrix}
		-\Delta_L & -2\Omega_c^* \\
		-2\Omega_c & \Delta_L
	\end{pmatrix}.
\end{equation}
If one defines the spectroscopic Rabi frequency by \(\Omega = 2|\Omega_c| = 2\mu_{eg}E_0\), the generalized Rabi frequency is
\begin{equation}
	\Omega_R = \sqrt{\Delta_L^2 + \Omega^2}.
	\label{eq:RabiFrequency}
\end{equation}
The phase \(\phi\) does not change \(\Omega_R\), but it rotates the drive axis in the Bloch sphere and therefore changes the phase of the oscillation.


\subsection{Recovery of Laboratory Frame Density Matrix}

To recover the laboratory-frame density matrix from the evolved rotating-frame density matrix, use
\begin{equation}
	\rho^{\mathrm{lab}}(t) = U^\dagger(t)\rho(t)U(t).
	\label{eq:RecoverOriginalDensityMatrix}
\end{equation}
The matrix elements are related by
\begin{align}
	\rho_{gg}^{\mathrm{lab}}(t) & = \rho_{gg}(t) \label{eq:RecoveryGG} \\
	\rho_{ee}^{\mathrm{lab}}(t) & = \rho_{ee}(t) \label{eq:RecoveryEE} \\
	\rho_{ge}^{\mathrm{lab}}(t) & = e^{i\omega_{L} t}\sigma_{ge}(t) \label{eq:RecoveryGE} \\
	\rho_{eg}^{\mathrm{lab}}(t) & = e^{-i\omega_{L} t}\sigma_{eg}(t). \label{eq:RecoveryEG}
\end{align}

%----------------------------------------------------------------------------------------
%	SECTION: Applications and Implications
%----------------------------------------------------------------------------------------

\section{Applications and Implications}

Rabi oscillations and the associated theoretical tools, such as the RWA and density matrix formalism, have wide-ranging applications:
\begin{itemize}
	\item \textbf{Quantum Computing:} Rabi oscillations are used to implement quantum gates by precisely controlling the population of qubits.
	\item \textbf{Spectroscopy:} The driven two-level model provides the basic reference for interpreting coherent optical signals.
	\item \textbf{Atomic Physics:} Understanding Rabi oscillations is essential for manipulating atomic states in experiments.
\end{itemize}
These concepts form the foundation for advanced topics in quantum mechanics and quantum technologies.
